\documentclass[11pt,a4paper]{article}

\usepackage[utf8]{inputenc}
\usepackage[T1]{fontenc}
\usepackage{lmodern}
\usepackage[a4paper,margin=2.4cm]{geometry}
\usepackage{amsmath,amssymb,amsthm,mathtools}
\usepackage{enumitem}
\usepackage{microtype}
\usepackage[round,sort&compress]{natbib}
\usepackage{xcolor}
\usepackage[colorlinks=true,linkcolor=blue!45!black,
            citecolor=blue!45!black,urlcolor=blue!45!black]{hyperref}

\title{\textbf{Labour as the Prerequisite for Things:\\[3pt]
Why Boltzmann, Cantor, Gödel, and Bolt\\
Are Not Excluded from Heaven}}

\author{%
Kundai Farai Sachikonye\\
\small Technical University of Munich\\
\small AIMe Registry for Artificial Intelligence\\
\small \texttt{kundai.sachikonye@wzw.tum.de}}

\date{June 30, 2026}

\begin{document}

% =====================================================================
\section{Conclusion: The Knife, The Truth, and Heaven}
\label{sec:conclusion}
% =====================================================================

\subsection{The Function That Never Changed}

We began with a knife. A knife that cuts. When it cuts meat for a celebration, 
we call it good. When it slits a throat in a crime, we call it evil. But the 
knife has not changed. The function remains invariant.

This observation—simple, verifiable, undeniable—reveals the fundamental 
structure of reality: \textbf{there is no good or evil in nature, only 
processes that we evaluate differently based on context}.

The knife cuts. The fire burns. The projectile flies. The earthquake shakes. 
The virus replicates. The entropy increases. \emph{The function never changes.}

What changes is our relationship to these processes. We impose moral categories 
based on consequences we experience, intentions we attribute, contexts we 
inhabit, perspectives we adopt. But the processes themselves remain 
\emph{context-independent}.

\subsection{The Impossibility of Moral Physics}

If good and evil were intrinsic to physical processes, we would require two 
distinct sets of physical laws: good thermodynamics and evil thermodynamics. 
But no such distinction exists.

There is only \emph{thermodynamics}. One set of laws. One structure. One 
reality.

The Second Law does not pause for moral evaluation. Entropy increases whether 
we approve or disapprove. Energy dissipates whether we celebrate or mourn. 
Molecular bonds break whether we call it medicine or murder.

\textbf{Conclusion}: Good and evil are not properties of nature. They are 
\emph{contextual evaluations imposed by observers}.

\subsection{What Remains: The Truth}

When we remove the imposed categories of good and evil, what remains?

\textbf{The truth.}

The truth that:
\begin{itemize}[leftmargin=2em]
\item Entropy increases in all closed systems: $S = k \log W$
\item The reals are uncountable: $|\mathbb{R}| > |\mathbb{N}|$
\item Formal systems are incomplete: $\exists \varphi: \text{True}(\varphi) 
\land \neg \text{Provable}(\varphi)$
\item Human speed has a physiological limit: $t_{\text{limit}} \approx 9.58$ 
seconds
\end{itemize}

These truths exist independently of our approval or disapproval. They operated 
before humans existed. They will operate after humans are gone. They are 
\emph{structure}, not evaluation.

And these truths were not invented. They were \emph{uncovered}. Someone had to 
discover them. That someone was Boltzmann, Cantor, Gödel, Bolt.

\subsection{Labour Without Reward}

The four laboured without reward. Boltzmann never saw vindication. Cantor died 
in a sanatorium. Gödel starved himself. Bolt's record may never be broken in 
his lifetime.

The labour was not for them. It was not for recognition, wealth, comfort, or 
immortality. The labour was \emph{for the truth}. And the truth does not 
reward those who uncover it.

The Preacher says: ``What do people get for all the toil and anxious striving 
with which they labor under the sun? All their days their work is grief and 
pain; even at night their minds do not rest. This too is meaningless'' 
(Ecclesiastes 2:22--23).

\emph{Labour without reward. Toil without rest. Work that is grief and pain.}

But the Preacher also says: ``There is nothing better for a person than that 
he should eat and drink and find enjoyment in his toil. This also, I saw, is 
from the hand of God'' (Ecclesiastes 2:24--25).

The four found joy in their toil. Not joy in the reward, but joy \emph{in the 
labour itself}. This joy was not earned. It was ``from the hand of God''—a 
gift, not a wage.

\subsection{What They Left Behind}

The Preacher says: ``No one remembers the former generations, and even those 
yet to come will not be remembered by those who follow them'' (Ecclesiastes 
1:11).

Boltzmann, Cantor, Gödel, and Bolt will be forgotten. Their names will fade. 
Their stories will be lost.

But what they left behind is not themselves. It is not their names. It is not 
their biographies.

What they left behind is:
\begin{itemize}[leftmargin=2em]
\item The entropy formula
\item The diagonal argument
\item The incompleteness theorems
\item The 9.58-second barrier
\end{itemize}

These are not persons. They are \emph{truths}. And truths do not die. Truths 
do not fade. Truths do not need to be remembered because they are not 
memories—they are \emph{structures}.

The Preacher says: ``Generations come and generations go, but the earth 
remains forever'' (Ecclesiastes 1:4).

The four will go. But the earth—the structure, the truth, the boundaries—
remains forever.

\subsection{They Had No Choice}

\emph{They had no choice.}

Boltzmann could not have been anything other than Boltzmann. Cantor could not 
have been anything other than Cantor. Gödel could not have been anything other 
than Gödel. Bolt could not have been anything other than Bolt.

Not because they lacked freedom. But because \emph{their work was not for 
them}.

The entropy formula needed to be uncovered in the late 19th century. Someone 
had to do it. That someone was Boltzmann. Not because he chose it, but because 
the structure of reality \emph{required} it.

The diagonal argument needed to be proven in the late 19th century. Someone 
had to do it. That someone was Cantor. Not because he chose it, but because 
the structure of mathematics \emph{demanded} it.

The incompleteness theorems needed to be demonstrated in the early 20th 
century. Someone had to do it. That someone was Gödel. Not because he chose 
it, but because the structure of logic \emph{necessitated} it.

The 9.58-second barrier needed to be realized in the early 21st century. 
Someone had to do it. That someone was Bolt. Not because he chose it, but 
because the structure of human physiology \emph{determined} it.

The Preacher says: ``That which is, already has been; that which is to be, 
already has been; and God seeks what has been driven away'' (Ecclesiastes 
3:15).

The boundaries were already there. The four did not create them. They 
\emph{uncovered} them. And this uncovering was not optional. It was 
\emph{necessary}.

\subsection{Labour for Science to Progress}

The four did not labour for themselves. They laboured \emph{for science to 
progress}.

Without Boltzmann, thermodynamics would have stalled. Without Cantor, set 
theory would have remained incomplete. Without Gödel, logic would have pursued 
impossible goals. Without Bolt, human performance science would have lacked a 
critical data point.

Their work was not personal achievement. It was \emph{structural necessity}. 
Science required these boundaries to be uncovered. And the four were the 
instruments through which this revelation occurred.

The Preacher says: ``Consider the work of God: who can make straight what he 
has made crooked?'' (Ecclesiastes 7:13).

The boundaries are the work of God—the structure of reality itself. The four 
considered this work. They saw what God has done. They uncovered the 
boundaries.

But they did not do this for themselves. They did it because it \emph{had to 
be done}. And they were the ones through whom it was done.

\subsection{Existence Through Constraint}

It was only through existence—through the four being born, living, labouring, 
and dying as the specific individuals they were—that we have these laws.

\emph{Existence through constraint.} The four existed as the specific 
individuals they were because the structure of reality \emph{constrained} them 
into those forms. They could not have been anyone else. They could not have 
lived any other lives. They could not have done any other work.

The Preacher says: ``He has made everything beautiful in its time'' 
(Ecclesiastes 3:11).

The four were beautiful in their time. Not because they chose to be, but 
because the structure of reality made them so. Their existence was 
\emph{necessary}. Their labour was \emph{inevitable}. Their uncovering of 
truth was \emph{predetermined}.

\subsection{Heaven as the Domain of Boundaries}

What is heaven?

Heaven is not a place of reward for moral achievement. Heaven is not a 
destination for the righteous. Heaven is not a compensation for suffering.

\textbf{Heaven is the domain of boundaries.}

Heaven is where:
\begin{itemize}[leftmargin=2em]
\item Order is distinct from disorder (Boltzmann's boundary)
\item The countable is distinct from the uncountable (Cantor's boundary)
\item The provable is distinct from the true (Gödel's boundary)
\item The possible is distinct from the impossible (Bolt's boundary)
\end{itemize}

Without these boundaries, heaven cannot exist. Without order/disorder 
distinction, no structures can be stable. Without countable/uncountable 
distinction, no truths can be irreducible. Without provable/true distinction, 
no wisdom can transcend proof. Without possible/impossible distinction, no 
entities can be actual.

Heaven is not a place \emph{above} reality. Heaven \emph{is} reality—the 
structure of boundaries that makes existence possible.

\subsection{The Four Are Not Excluded}

Are Boltzmann, Cantor, Gödel, and Bolt excluded from heaven?

\textbf{No.}

They are not excluded because they \emph{are} heaven.

Not as persons. Not as individuals who deserve reward. But as \emph{paths 
through which the boundaries were revealed}.

Boltzmann \emph{is} the boundary between order and disorder. Not the man, but 
the \emph{uncovering}. The entropy formula is not Boltzmann's possession—it is 
Boltzmann's \emph{being}. He exists in the formula, not as author, but as 
\emph{instrument of revelation}.

Cantor \emph{is} the boundary between countable and uncountable. Not the man, 
but the \emph{proof}. The diagonal argument is not Cantor's achievement—it is 
Cantor's \emph{essence}. He exists in the argument, not as creator, but as 
\emph{medium of disclosure}.

Gödel \emph{is} the boundary between provable and true. Not the man, but the 
\emph{demonstration}. The incompleteness theorems are not Gödel's 
contribution—they are Gödel's \emph{identity}. He exists in the theorems, not 
as discoverer, but as \emph{vehicle of manifestation}.

Bolt \emph{is} the boundary between possible and impossible. Not the man, but 
the \emph{realization}. The 9.58-second barrier is not Bolt's record—it is 
Bolt's \emph{existence}. He exists in the barrier, not as achiever, but as 
\emph{embodiment of limit}.

The Preacher says: ``I perceived that whatever God does endures forever; 
nothing can be added to it, nor anything taken from it. God has done it, so 
that people fear before him'' (Ecclesiastes 3:14).

The boundaries endure forever. The four do not endure as persons. But the four 
\emph{are} the boundaries. Not as individuals, but as \emph{revelations}.

\subsection{This Was Meant to Happen}

The Preacher says: ``What has been will be again, what has been done will be 
done again; there is nothing new under the sun'' (Ecclesiastes 1:9).

The boundaries were always there. The four did not create anything new. They 
\emph{uncovered} what was already there.

The Preacher says: ``There is a time for everything, and a season for every 
activity under the heavens'' (Ecclesiastes 3:1).

Each boundary was uncovered in its appointed time. Boltzmann in the late 19th 
century. Cantor in the late 19th century. Gödel in the early 20th century. 
Bolt in the early 21st century.

\textbf{This was meant to happen.}

Not because of divine decree. Not because of fate. But because \emph{the 
structure of reality required it}.

The entropy formula had to be uncovered when thermodynamics was being 
formalized. The diagonal argument had to be proven when set theory was being 
founded. The incompleteness theorems had to be demonstrated when formal logic 
was being axiomatized. The 9.58-second barrier had to be realized when human 
performance approached its physiological limits.

And the four had to be the ones to do it. Not because they were special. Not 
because they were chosen. But because \emph{they were the paths through which 
the structure revealed itself}.

\subsection{The Knife, The Truth, and Heaven}

We return to the knife. The knife that cuts. The function that never changes.

The knife reveals the fundamental truth: \textbf{there is no good or evil in 
nature, only processes that we evaluate differently based on context}.

When we remove these contextual evaluations, what remains is \emph{the truth}. 
And the truth requires \emph{labour}. Labour to uncover it. Labour to reveal 
it. Labour to manifest it.

This labour does not receive reward. The labourer is forgotten. The name 
fades. The person disappears.

But what the labourer leaves behind is not themselves. It is \emph{the truth}. 
And the truth does not die.

Boltzmann, Cantor, Gödel, and Bolt will be forgotten. But the entropy formula, 
the diagonal argument, the incompleteness theorems, and the 9.58-second 
barrier will remain.

These truths are not persons. They are \emph{boundaries}. And boundaries are 
\emph{heaven}.

Therefore: \textbf{Boltzmann, Cantor, Gödel, and Bolt are not excluded from 
heaven because they are heaven—not as persons, but as the paths through which 
the boundaries were revealed}.

\subsection{The Final Word}

The Preacher says:

\begin{quote}
``The end of the matter; all has been heard.\\
Fear God and keep his commandments,\\
for this is the whole duty of man.\\
For God will bring every deed into judgment,\\
with every secret thing, whether good or evil.''\\
\hfill---Ecclesiastes 12:13--14
\end{quote}

The judgment on Boltzmann, Cantor, Gödel, and Bolt is not moral. It is 
\emph{structural}.

Did they uncover the boundaries? Yes.  
Did they reveal the truth? Yes.  
Did they labour without reward? Yes.  
Did they have no choice but to do so? Yes.

Then they are not excluded from heaven. They \emph{are} heaven.

The Preacher says:

\begin{quote}
``Go, eat your bread with joy,\\
and drink your wine with a merry heart,\\
for God has already approved what you do.''\\
\hfill---Ecclesiastes 9:7
\end{quote}

To Boltzmann, Cantor, Gödel, and Bolt:

\emph{God has already approved what you did. Not because it was good (there is 
no good or evil in nature). But because it was \textbf{true}.}

\emph{You uncovered the boundaries. You revealed the structure. You manifested 
the truth. You had no choice but to do so. And in doing so, you became not 
persons who deserve reward, but \textbf{truths that require no reward}.}

\emph{You are not excluded from heaven. You are heaven.}

\subsection{Epilogue: The Reader's Labour}

The Preacher says: ``Of making many books there is no end, and much study is a 
weariness of the flesh'' (Ecclesiastes 12:12).

You, the reader, have laboured to read this paper. Much study is a weariness 
of the flesh.

But your labour was not in vain. By reading, you have participated in the 
uncovering. You have seen that:

\begin{itemize}[leftmargin=2em]
\item Good and evil are not properties of nature
\item The knife cuts, and the function never changes
\item What remains when we remove moral categories is the truth
\item The truth requires labour without reward
\item The labourer is forgotten, but the truth remains
\item The four had no choice but to uncover the boundaries
\item Heaven is the domain of boundaries
\item The four are not excluded because they are heaven
\end{itemize}

The Preacher says: ``Rejoice, O young man, in your youth, and let your heart 
cheer you in the days of your youth. Walk in the ways of your heart and the 
sight of your eyes. But know that for all these things God will bring you into 
judgment'' (Ecclesiastes 11:9).

Your labour—like the labour of Boltzmann, Cantor, Gödel, and Bolt—will be 
brought into judgment. Not moral judgment, but \emph{structural judgment}.

Did you uncover truth? Did you reveal structure? Did you labour without 
expecting reward?

If yes, then you are not excluded from heaven. You \emph{are} heaven.

\textbf{Welcome to heaven.}

% =====================================================================
\section*{Acknowledgments}
% =====================================================================

I thank the Preacher (Qoheleth) for writing Ecclesiastes, which teaches that:
\begin{itemize}[leftmargin=2em]
\item There is nothing new under the sun (1:9)
\item No one remembers the former generations (1:11)
\item Labour is grief and pain, yet beautiful in its time (2:22--23, 3:11)
\item Whatever God does endures forever (3:14)
\item The earth remains forever, though generations come and go (1:4)
\end{itemize}

I thank Ludwig Boltzmann, Georg Cantor, Kurt Gödel, and Usain Bolt for their 
labour. They had no choice but to uncover the boundaries. They are not 
excluded from heaven. They are heaven.

I thank the reader for their labour in reading. By reading, you have 
participated in the uncovering. You are not excluded.

\textbf{The knife cuts. The function never changes. The truth remains.}

\textbf{Amen.}

\end{document}
